\chapter[Resultados Experimentais]{Resultados Experimentais}
\label{cap:resultados}

Este capítulo apresenta os resultados experimentais obtidos no benchmark Families in the Wild (FIW). A comparação principal segue o protocolo descrito no Capítulo~\ref{cap:Metodologia}: validação cruzada de cinco dobras separadas por família, escolha do limiar apenas na validação e avaliação no conjunto de teste do RFIW Track-I. Entram nessa comparação a linha de base AdaFace-Cosseno, o ViT-FaCoR, o DINOv2-LoRA, o Modelo com Recuperação e as três configurações finalistas do AdaFace-Regional: Simétrica, Comparativa e com Negativos Difíceis. Os VLMs são apresentados em seção própria como linhas de base externas zero-shot, avaliadas em amostras balanceadas de 6.000 pares e sem validação cruzada.

A métrica principal é o ROC AUC, por ser independente do limiar de decisão. Average Precision e TAR@FAR complementam essa leitura ao avaliar, respectivamente, a curva precisão-revocação e os regimes de baixa falsa aceitação. A acurácia é usada principalmente na análise por tipo de relação, porque permite observar quais vínculos familiares continuam mais difíceis para cada configuração.

\section{Resultados gerais em FIW}

A Tabela~\ref{tab:fiw_main} resume os resultados dos modelos avaliados sob validação cruzada em FIW. Todos os modelos supervisionados da comparação principal, incluindo DINOv2-LoRA e Modelo com Recuperação, foram avaliados no regime de cinco dobras. Para o AdaFace-Regional, são mantidas três configurações finalistas porque elas representam escolhas operacionais próximas, mas não idênticas.

A configuração Simétrica consolida o ganho principal do modelo, a Comparativa desloca o ponto de operação para baixa falsa aceitação e a configuração com Negativos Difíceis combina essa estratégia com pares negativos mais próximos dos positivos. AdaFace-Cosseno, DINOv2-LoRA e Modelo com Recuperação aparecem na mesma tabela para situar, respectivamente, a linha de base sem treinamento, a adaptação com DINOv2 e a alternativa baseada em recuperação de pares positivos. Os VLMs não entram nessa tabela, porque foram avaliados em outro cenário experimental.

\begin{table}[htbp]
\centering
\scriptsize
\setlength{\tabcolsep}{3pt}
\caption{Resultados gerais em FIW sob validação cruzada de cinco dobras. Os valores reportam média $\pm$ desvio padrão; no AdaFace-Cosseno, o desvio é zero porque o extrator é congelado e o conjunto de teste é fixo.}
\label{tab:fiw_main}
\begin{tabularx}{\textwidth}{Xccccc}
\hline
\textbf{Modelo} & \textbf{AUC} & \textbf{AP} & \textbf{TAR@0{,}001} & \textbf{TAR@0{,}01} & \textbf{TAR@0{,}1} \\
\hline
AdaFace-Regional com Negativos Difíceis & \textbf{0{,}876 $\pm$ 0{,}003} & \textbf{0{,}856 $\pm$ 0{,}003} & 0{,}068 $\pm$ 0{,}015 & 0{,}207 $\pm$ 0{,}011 & \textbf{0{,}595 $\pm$ 0{,}008} \\
AdaFace-Regional Comparativo & 0{,}874 $\pm$ 0{,}004 & 0{,}853 $\pm$ 0{,}005 & 0{,}059 $\pm$ 0{,}012 & 0{,}192 $\pm$ 0{,}014 & 0{,}589 $\pm$ 0{,}015 \\
AdaFace-Regional Simétrico & 0{,}873 $\pm$ 0{,}004 & 0{,}848 $\pm$ 0{,}004 & 0{,}042 $\pm$ 0{,}013 & 0{,}174 $\pm$ 0{,}012 & 0{,}577 $\pm$ 0{,}009 \\
ViT-FaCoR & 0{,}846 $\pm$ 0{,}004 & 0{,}813 $\pm$ 0{,}005 & 0{,}022 $\pm$ 0{,}006 & 0{,}135 $\pm$ 0{,}012 & 0{,}496 $\pm$ 0{,}026 \\
DINOv2-LoRA com Descongelamento Completo & 0{,}814 $\pm$ 0{,}007 & 0{,}784 $\pm$ 0{,}009 & 0{,}013 $\pm$ 0{,}002 & 0{,}102 $\pm$ 0{,}009 & 0{,}469 $\pm$ 0{,}018 \\
AdaFace-Cosseno, sem treino & 0{,}799 $\pm$ 0{,}000 & 0{,}809 $\pm$ 0{,}000 & \textbf{0{,}071 $\pm$ 0{,}000} & \textbf{0{,}218 $\pm$ 0{,}000} & 0{,}524 $\pm$ 0{,}000 \\
Modelo com Recuperação & 0{,}778 $\pm$ 0{,}008 & 0{,}741 $\pm$ 0{,}015 & 0{,}014 $\pm$ 0{,}002 & 0{,}078 $\pm$ 0{,}019 & 0{,}384 $\pm$ 0{,}030 \\
\hline
\end{tabularx}
\end{table}

O maior valor de AUC na comparação principal em validação cruzada é obtido pelo AdaFace-Regional com Negativos Difíceis, com 0,876 $\pm$ 0{,}003. Esse valor corresponde à média e ao desvio padrão das cinco dobras, sem agrupar predições de modelos treinados em partições distintas. Assim, ele é o resultado principal usado para representar o desempenho da configuração final no protocolo adotado.

A comparação mais importante é entre o AdaFace-Regional e o ViT-FaCoR, porque o ViT-FaCoR representa a principal arquitetura supervisionada anterior do estudo. A configuração com Negativos Difíceis alcança AUC médio de 0,876, contra 0,846 do ViT-FaCoR, diferença de 0,030 no valor de AUC. As outras configurações finalistas do AdaFace-Regional apresentam leitura semelhante: a Simétrica e a Comparativa também ficam cerca de 0,027 acima no valor de AUC em relação ao ViT-FaCoR. Como o desvio por dobra está em torno de 0,004 nesses modelos, a diferença observada é maior do que a variação média entre dobras. Esse é o principal resultado experimental do trabalho: o conjunto de componentes formado por backbone facial especializado, tokens regionais anatômicos, cabeça auxiliar de relação e passagem simétrica está associado a ganho empírico sobre o ViT-FaCoR no FIW sob o protocolo adotado.

A análise pareada das cinco dobras reforça essa interpretação para a métrica principal. Comparando AdaFace-Regional com Negativos Difíceis e ViT-FaCoR, a diferença média é de +0,0299 em ROC AUC, e o teste t pareado resultou em $t(4)=10{,}264$ e $p=0{,}0005$. As configurações Simétrica e Comparativa também apresentam AUC superior à do ViT-FaCoR no teste t pareado, com $p=0{,}0014$ e $p=0{,}0013$, respectivamente. No teste de Wilcoxon signed-rank bilateral, entretanto, as três comparações resultam em $p=0{,}0625$, valor influenciado pelo número reduzido de observações ($n=5$ dobras). Por isso, os testes são lidos como evidência complementar à média e ao desvio por dobra, e não como substituto da análise operacional por métrica.

A escolha entre as três configurações finalistas exige mais cuidado. Em AUC, elas ficam muito próximas: 0,873, 0,874 e 0,876, todas dentro de uma faixa menor ou comparável ao desvio entre dobras. Essa sobreposição não sustenta uma afirmação forte de superioridade entre as três variantes; ela indica, antes, que as configurações ocupam uma região de desempenho semelhante. A diferença entre elas aparece mais claramente no ponto de operação. Em TAR@FAR=0{,}001, a configuração com Negativos Difíceis chega a 0,068 $\pm$ 0,015, enquanto a Simétrica fica em 0,042 $\pm$ 0,013. Assim, as três configurações são próximas em ranking global, mas, entre as variantes do AdaFace-Regional, a configuração com Negativos Difíceis é a escolha mais adequada quando falsos positivos têm custo elevado.

O comportamento do AdaFace-Cosseno reforça a importância de separar AUC global e regime operacional. Mesmo sem treinamento específico para parentesco, essa linha de base permanece forte nas faixas de baixa falsa aceitação. Em TAR@FAR=0{,}01, o valor do AdaFace-Cosseno (0,218 $\pm$ 0,000) fica no mesmo patamar da configuração com Negativos Difíceis (0,207 $\pm$ 0,011). Em TAR@FAR=0{,}001, AdaFace-Cosseno (0,071 $\pm$ 0,000) fica ligeiramente acima dos 0,068 $\pm$ 0,015 da configuração com Negativos Difíceis. Isso indica que embeddings de reconhecimento facial já carregam um sinal útil para parentesco, principalmente quando a comparação exige baixa taxa de falso positivo. A contribuição do AdaFace-Regional, portanto, não é substituir esse sinal em todos os pontos de operação, mas melhorar a ordenação global e oferecer um modelo adaptado à tarefa.

Sob validação cruzada, o DINOv2-LoRA alcança AUC de 0,814 $\pm$ 0,007, ficando acima da linha de base AdaFace-Cosseno em AUC, mas abaixo do ViT-FaCoR. O resultado indica que a adaptação do DINOv2 por LoRA e descongelamento completo melhora a linha de base sem treinamento na ordenação global, porém não alcança a arquitetura com ViT e atenção cruzada bidirecional. O Modelo com Recuperação apresenta AUC de 0,778 $\pm$ 0,008, abaixo do AdaFace-Cosseno, e permanece como hipótese útil para investigar se exemplos positivos recuperados podem reduzir diferenças entre classes frequentes e raras. Dessa forma, DINOv2-LoRA e Modelo com Recuperação ajudam a delimitar duas leituras negativas: a representação auto-supervisionada adaptada não resolve sozinha o problema, e a recuperação de exemplos não melhora necessariamente o ranking global.

Os VLMs binários, avaliados em protocolo separado, mostram um limite claro da inferência zero-shot. O Claude Sonnet atinge AUC de 0,789, muito próximo do AdaFace-Cosseno sem treinamento específico (0,799), mas ainda abaixo dele. O Codex VLM fica mais distante, com AUC de 0,624 e Average Precision de 0,592. A diferença principal aparece nos regimes de baixa falsa aceitação: o Claude mantém TAR@FAR=0{,}01 de apenas 0,042, e o Codex chega a 0,018. Mesmo quando a decisão binária produz acurácia ou F1 razoáveis, a ordenação dos escores não é fina o suficiente para operar em pontos conservadores de falsa aceitação.

\section{Execuções do DINOv2-LoRA e do AdaFace-Regional}

Esta seção registra as execuções de desenvolvimento do DINOv2-LoRA e do AdaFace-Regional. O objetivo não é substituir a comparação principal da Tabela~\ref{tab:fiw_main}, mas mostrar como as configurações finais foram escolhidas. No DINOv2-LoRA, as execuções de desenvolvimento serviram para selecionar a configuração com descongelamento completo. No AdaFace-Regional, o ciclo foi mais longo, porque nele está a proposta autoral do trabalho.

O DINOv2-LoRA foi avaliado em sete execuções, cobrindo regimes de adaptação e descongelamento do DINOv2. A Tabela~\ref{tab:dinov2_lora_execucoes} mostra que o melhor AUC da família aparece no descongelamento completo, enquanto o melhor TAR@FAR=0{,}01 ocorre na configuração inicial com LoRA. Isso impede uma leitura simples de que liberar mais parâmetros sempre melhora todos os regimes de avaliação.

\begin{table}[htbp]
\centering
\scriptsize
\setlength{\tabcolsep}{3pt}
\caption{Execuções de desenvolvimento do DINOv2-LoRA em FIW.}
\label{tab:dinov2_lora_execucoes}
\begin{tabularx}{\textwidth}{Xrrrrr}
\hline
\textbf{Regime} & \textbf{Params} & \textbf{AUC} & \textbf{Acc.} & \textbf{TAR@0{,}01} & \textbf{Gap val-teste} \\
\hline
LoRA-only                     & 8{,}5M  & 0{,}806 & 72{,}6\% & \textbf{0{,}152} & -0{,}105 \\
regularização mais forte      & 8{,}5M  & 0{,}799 & 72{,}4\% & 0{,}095          & -0{,}106 \\
sampler estratificado         & 8{,}5M  & 0{,}809 & 71{,}0\% & 0{,}098          & -0{,}100 \\
descongelamento parcial       & 36{,}8M & 0{,}812 & 72{,}9\% & 0{,}119          & -0{,}099 \\
descongelamento completo      & 93{,}5M & \textbf{0{,}822} & 72{,}0\% & 0{,}100 & \textbf{-0{,}071} \\
perda contrastiva dominante   & 8{,}5M  & 0{,}814 & 72{,}9\% & 0{,}094          & -0{,}100 \\
DINOv2 + ViT-FaCoR            & 11{,}6M & 0{,}810 & 71{,}9\% & 0{,}136          & -0{,}097 \\
\hline
\end{tabularx}
\end{table}

A leitura do DINOv2-LoRA é desfavorável à hipótese de que DINOv2 congelado com LoRA e atenção diferencial seria suficiente para alcançar o ViT-FaCoR. Entre as execuções de desenvolvimento, o melhor AUC global aparece no descongelamento completo, enquanto a configuração inicial com LoRA preserva o melhor TAR@FAR=0{,}01. Esse comportamento mostra que o ganho em AUC ocorre em detrimento do regime de baixa falsa aceitação. A variante híbrida com pesos do ViT-FaCoR também não transfere o desempenho do sistema original, o que sugere que o resultado do ViT-FaCoR depende do conjunto formado por backbone, atenção cruzada, função de perda e calibração, e não apenas do encoder isolado.

O AdaFace-Regional teve o ciclo de desenvolvimento mais extenso. A Tabela~\ref{tab:adaface_regional_execucoes} resume as principais execuções em FIW. O descongelamento parcial é o primeiro salto relevante, porque libera parte do AdaFace. A configuração Simétrica introduz a passagem nas duas ordens e consolida o ganho central da arquitetura. As configurações seguintes ajustam o ponto de operação, com destaque para a configuração com Negativos Difíceis, que acrescenta pares negativos mais difíceis ao treino.

\begin{table}[htbp]
\centering
\scriptsize
\setlength{\tabcolsep}{3pt}
\caption{Execuções de desenvolvimento do AdaFace-Regional em FIW.}
\label{tab:adaface_regional_execucoes}
\begin{tabularx}{\textwidth}{Xrrrrr}
\hline
\textbf{Intervenção principal} & \textbf{AUC} & \textbf{AP} & \textbf{TAR@0{,}001} & \textbf{TAR@0{,}01} & \textbf{Gap val-teste} \\
\hline
AdaFace congelado               & 0{,}746 & 0{,}732 & 0{,}024 & 0{,}101 & -0{,}089 \\
descongelamento parcial          & 0{,}856 & 0{,}839 & 0{,}042 & 0{,}176 & -0{,}076 \\
descongelamento parcial + perda contrastiva & 0{,}851 & 0{,}831 & 0{,}027 & 0{,}166 & -0{,}080 \\
reamostragem preliminar de negativos & 0{,}847 & 0{,}829 & 0{,}020 & 0{,}147 & -0{,}088 \\
cabeça auxiliar de relação       & 0{,}848 & 0{,}829 & 0{,}029 & 0{,}167 & -0{,}084 \\
passagem simétrica               & 0{,}879 & 0{,}856 & 0{,}033 & 0{,}186 & \textbf{-0{,}026} \\
taxas de aprendizado diferenciadas & 0{,}873 & 0{,}852 & 0{,}045 & 0{,}196 & -0{,}036 \\
regularização L2-SP              & 0{,}879 & 0{,}856 & 0{,}028 & 0{,}191 & \textbf{-0{,}026} \\
fusão apenas comparativa         & 0{,}874 & 0{,}850 & 0{,}054 & 0{,}178 & -0{,}027 \\
fusão comparativa + taxas diferenciadas & 0{,}875 & 0{,}857 & 0{,}052 & 0{,}212 & -0{,}030 \\
negativos difíceis reais (30\%)  & \textbf{0{,}882} & \textbf{0{,}865} & \textbf{0{,}075} & \textbf{0{,}213} & \textbf{-0{,}021} \\
\hline
\end{tabularx}
\end{table}

No ciclo de desenvolvimento, duas mudanças aparecem associadas aos maiores saltos de desempenho. A primeira é o descongelamento parcial: com AdaFace congelado, o modelo fica em AUC de 0,746, enquanto o descongelamento parcial sobe para 0,856 ao liberar o último estágio e a camada de saída. A segunda é a passagem simétrica, que processa cada par nas duas ordens, $(A,B)$ e $(B,A)$, e reduz a dependência da posição das imagens. Essa mudança aumentou o AUC no ciclo de desenvolvimento em 0,022 em relação ao descongelamento parcial e reduziu o gap validação-teste de -0,076 para -0,026.

A configuração com Negativos Difíceis fecha o ciclo ao avaliar pares negativos com a mesma estrutura relacional anotada dos positivos. Nessa configuração, 30\% dos negativos de treino são formados por pessoas de famílias diferentes, mas com os mesmos papéis do par positivo de referência. Esses pares aumentam a exigência discriminativa da cabeça de decisão sem violar a separação por família. O efeito aparece principalmente na especificidade de não-parentes e nas métricas de baixa falsa aceitação, ainda que com custo em algumas classes positivas.

Também é importante registrar as tentativas que não produziram ganho claro. A perda contrastiva auxiliar não melhorou o AUC e prejudicou classes de avós e netos. A reamostragem preliminar de negativos ficou abaixo do descongelamento parcial e não foi usada para sustentar a análise final. A regularização L2-SP empatou numericamente com a passagem simétrica, sem ganho mensurável. As taxas de aprendizado diferenciadas, isoladamente, também ficaram dentro do ruído. Assim, as intervenções com evidência mais consistente no ciclo foram o descongelamento parcial, a passagem simétrica e os negativos difíceis reais.

\section{Análise por tipo de relação}

A análise por tipo de relação permite verificar se o ganho do AdaFace-Regional é distribuído entre os vínculos ou concentrado em alguns grupos. A Tabela~\ref{tab:per_rel_compare} apresenta a acurácia por relação para as três configurações finalistas do AdaFace-Regional, como média e desvio padrão sobre cinco dobras.

\begin{table}[htbp]
\centering
\small
\setlength{\tabcolsep}{3pt}
\caption{Acurácia (\%) por tipo de relação em FIW, como média $\pm$ desvio padrão sobre cinco dobras, para as três configurações finalistas do AdaFace-Regional. O N de pares é o mesmo em todas as dobras, porque o conjunto de teste do RFIW Track-I é fixo.}
\label{tab:per_rel_compare}
\begin{tabular}{lrrrr}
\hline
\textbf{Relação} & \textbf{N pares} & \textbf{Simétrica} & \textbf{Comparativa} & \textbf{Neg. difíceis} \\
\hline
fs (pai-filho)      & 1.135 & \textbf{90{,}7 $\pm$ 1{,}7} & 89{,}4 $\pm$ 1{,}3 & 89{,}5 $\pm$ 1{,}3 \\
fd (pai-filha)      & 918   & \textbf{86{,}9 $\pm$ 2{,}5} & 85{,}1 $\pm$ 3{,}3 & 84{,}7 $\pm$ 1{,}4 \\
ms (mãe-filho)      & 1.036 & 87{,}4 $\pm$ 3{,}4 & 86{,}2 $\pm$ 3{,}8 & \textbf{88{,}2 $\pm$ 2{,}6} \\
md (mãe-filha)      & 1.038 & \textbf{91{,}9 $\pm$ 2{,}0} & 90{,}4 $\pm$ 3{,}8 & 91{,}0 $\pm$ 1{,}1 \\
bb (irmão-irmão)    & 860   & \textbf{91{,}9 $\pm$ 1{,}4} & 89{,}8 $\pm$ 2{,}6 & 90{,}9 $\pm$ 1{,}4 \\
ss (irmã-irmã)      & 731   & \textbf{92{,}2 $\pm$ 2{,}0} & 89{,}2 $\pm$ 2{,}5 & 89{,}7 $\pm$ 1{,}4 \\
sibs (irmãos mistos) & 234   & \textbf{94{,}1 $\pm$ 1{,}1} & 91{,}6 $\pm$ 2{,}4 & 92{,}7 $\pm$ 2{,}3 \\
\hline
gfgs (avô-neto)     & 98  & \textbf{83{,}5 $\pm$ 2{,}3} & 82{,}0 $\pm$ 5{,}6 & 78{,}8 $\pm$ 3{,}1 \\
gfgd (avô-neta)     & 138 & \textbf{84{,}3 $\pm$ 2{,}1} & 84{,}2 $\pm$ 3{,}0 & 83{,}3 $\pm$ 1{,}5 \\
gmgs (avó-neto)     & 121 & \textbf{80{,}0 $\pm$ 2{,}8} & 73{,}6 $\pm$ 5{,}0 & 71{,}1 $\pm$ 3{,}8 \\
gmgd (avó-neta)     & 123 & \textbf{83{,}7 $\pm$ 5{,}6} & 80{,}8 $\pm$ 6{,}9 & 74{,}8 $\pm$ 4{,}7 \\
\hline
non-kin             & 6.993 & 67{,}5 $\pm$ 2{,}8 & \textbf{69{,}7 $\pm$ 4{,}4} & \textbf{69{,}7 $\pm$ 1{,}8} \\
\hline
\end{tabular}
\end{table}

Nas relações entre pais e filhos, as três configurações ficam em patamar alto, entre 85\% e 92\%. A configuração Simétrica apresenta os maiores valores em pai-filho, pai-filha e mãe-filha, enquanto a configuração com Negativos Difíceis tem o maior valor em mãe-filho. O desvio entre dobras é relativamente pequeno nessas classes, porque elas possuem mais pares no conjunto de teste. O ganho sobre o descongelamento parcial sugere que o conjunto de componentes finalistas está associado à redução de atalhos ligados à posição das imagens no par.

Nas relações entre irmãos, o desempenho é ainda mais estável. As classes bb, ss e sibs ficam próximas ou acima de 90\%, com desvios menores do que os observados nas relações de avós e netos. Uma explicação possível é que irmãos tendem a apresentar menor diferença geracional do que a observada em pares de pais e filhos ou de avós e netos, o que reduz parte da variação visual associada à idade.

O avanço mais relevante aparece nas relações de avós e netos. No descongelamento parcial, essas classes eram o principal gargalo do modelo, com acurácias entre 37\% e 52\%. Sob validação cruzada, a configuração Simétrica leva essas relações para a faixa de 80\% a 84\%. Esse ganho, porém, deve ser lido com cautela, porque são as classes com menos pares no teste, variando de 98 a 138 pares. Por isso, o desvio entre dobras é maior, especialmente em gmgd e gmgs. As configurações Comparativa e com Negativos Difíceis ficam um pouco abaixo da Simétrica nessas relações, principalmente em avó-neto e avó-neta, o que mostra o custo de deslocar o modelo para maior rejeição de falsos positivos.

A classe non-kin apresenta o movimento oposto. A configuração Simétrica é a mais forte nas classes positivas, mas tem menor acurácia em não-parentes, com 67,5 $\pm$ 2,8. As configurações Comparativa e com Negativos Difíceis sobem para 69,7, indicando maior especificidade. Esse comportamento é coerente com as métricas TAR@FAR: configurações que rejeitam melhor pares não aparentados tendem a perder parte da acurácia em relações positivas. A escolha entre esses regimes depende do custo relativo entre falso positivo e falso negativo na aplicação.

\section{Linhas de base com modelos multimodais}

Os VLMs foram avaliados novamente na formulação binária kin/non-kin, para alinhar a pergunta com a dos modelos supervisionados. Cada modelo recebeu 6.000 pares balanceados do FIW, com 3.000 positivos e 3.000 negativos, em regime zero-shot. A Tabela~\ref{tab:vlm_binario_resultados} resume as métricas globais, usando o escore contínuo derivado da decisão e da confiança conforme a Equação~\ref{eq:vlm_score}.

\begin{table}[htbp]
\centering
\caption{Resultados dos VLMs em verificação binária zero-shot no FIW, com 6.000 pares balanceados por modelo.}
\label{tab:vlm_binario_resultados}
\begin{tabularx}{0.82\textwidth}{Xcc}
\hline
\textbf{Métrica} & \textbf{Claude Sonnet} & \textbf{Codex VLM} \\
\hline
AUC & 0{,}789 & 0{,}624 \\
Average Precision & 0{,}742 & 0{,}592 \\
TAR@0{,}01 & 0{,}042 & 0{,}018 \\
Acurácia & 72{,}3\% & 59{,}0\% \\
Precisão & 0{,}691 & 0{,}563 \\
Revocação & 0{,}806 & 0{,}799 \\
Especificidade & 0{,}639 & 0{,}381 \\
F1 & 0{,}744 & 0{,}661 \\
Confiança média & 0{,}656 & 0{,}721 \\
\hline
\end{tabularx}
\end{table}

O Claude Sonnet apresenta desempenho acima do acaso e fica próximo da linha de base AdaFace-Cosseno em AUC, mas ligeiramente abaixo dela: 0,789 contra 0,799. A acurácia de 72,3\% e o F1 de 0,744 mostram que o modelo consegue responder à tarefa binária de forma coerente em muitos casos, porém a Average Precision e o TAR@FAR=0{,}01 indicam limitação forte de ordenação. A distribuição de confiança é discreta, com poucos valores recorrentes, e isso reduz a utilidade do modelo em regimes de falsa aceitação estrita.

O Codex VLM tem comportamento mais enviesado para a classe positiva. Sua revocação para kin é 0,799, parecida com a do Claude, mas a especificidade cai para 0,381. Em termos absolutos, isso significa que 1.857 dos 3.000 pares non-kin foram classificados como parentes. A matriz de confusão da Tabela~\ref{tab:vlm_confusao_binaria} explicita esse padrão. O resultado sugere que, em zero-shot, o Codex apresenta viés de aceitação positiva, com baixa especificidade diante de pares não biológicos com semelhança visual fortuita.

\begin{table}[htbp]
\centering
\small
\setlength{\tabcolsep}{5pt}
\caption{Matriz de confusão dos VLMs binários em 6.000 pares balanceados.}
\label{tab:vlm_confusao_binaria}
\begin{tabular}{lrrrr}
\hline
\textbf{Modelo} & \textbf{TP} & \textbf{FN} & \textbf{FP} & \textbf{TN} \\
\hline
Claude Sonnet 4.6 & 2.419 & 581 & 1.082 & 1.918 \\
Codex gpt-5.4-mini & 2.396 & 604 & 1.857 & 1.143 \\
\hline
\end{tabular}
\end{table}

A comparação com os modelos supervisionados deve ser feita com cautela. Os VLMs respondem à mesma tarefa binária, mas não têm treinamento no FIW, não usam validação cruzada e não ajustam limiar por validação. Ainda assim, as métricas independentes de limiar são informativas: o Claude fica próximo da linha de base facial sem treinamento, mas abaixo dela; o Codex fica substancialmente abaixo; e ambos ficam distantes do AdaFace-Regional e do ViT-FaCoR. Isso indica que conhecimento visual geral não substitui adaptação supervisionada para verificação de parentesco facial.

A avaliação multiclasse anterior permanece como diagnóstico histórico. Nela, o Codex zero-shot alcançou 33,1\% de acurácia e macro-F1 de 0,257 em 750 pares positivos; a adaptação por prompt caiu para 26,7\% de acurácia e macro-F1 de 0,223; e o piloto do Claude Sonnet em 75 pares ficou em torno de 37\% de acurácia. Como esse protocolo não inclui non-kin, ele não é usado na comparação principal. O recorte de 12 pares de avós e netos com prompts direcionados mostrou que orientar explicitamente o modelo para relações de duas gerações ajuda a escolher o grupo correto, mas quase não resolve a identificação da relação exata.

\section{Síntese dos resultados}

Sob validação cruzada, o AdaFace-Regional obteve o maior AUC entre as arquiteturas avaliadas em FIW. O resultado da configuração com Negativos Difíceis é de 0,876 $\pm$ 0,003 em AUC, acima do ViT-FaCoR (0,846 $\pm$ 0,004), do DINOv2-LoRA (0,814 $\pm$ 0,007), da linha de base AdaFace-Cosseno (0,799 $\pm$ 0,000) e do Modelo com Recuperação (0,778 $\pm$ 0,008). A principal evidência do capítulo é, portanto, que o conjunto de componentes da proposta está associado a melhor desempenho no FIW em relação às alternativas avaliadas dentro do protocolo adotado.

Fora da comparação principal em validação cruzada, os VLMs foram avaliados como linhas de base externas zero-shot. Nesse cenário, o Claude Sonnet binário alcança AUC de 0,789 e o Codex VLM binário alcança AUC de 0,624, ambos sem validação cruzada.

Essa conclusão deve ser acompanhada de quatro ressalvas. Primeiro, AdaFace-Cosseno continua competitivo nos regimes de baixa falsa aceitação, mostrando que reconhecimento facial pré-treinado já contém sinal útil para parentesco. Segundo, as diferenças entre as configurações Simétrica, Comparativa e com Negativos Difíceis são pequenas em AUC, e a escolha entre elas depende mais do ponto operacional do que de uma superioridade global clara. Terceiro, as relações de avós e netos continuam sendo as mais sensíveis à variação entre dobras, apesar do ganho expressivo obtido pelo AdaFace-Regional. Quarto, os VLMs binários zero-shot ficam abaixo dos modelos supervisionados e apresentam limitação clara em especificidade ou em operação com FAR baixo, indicando que inferência multimodal genérica não substitui modelos treinados e calibrados para esse domínio.
